\documentclass[12pt,a4paper]{article}
\usepackage[UTF8]{ctex}
\usepackage{amsmath,amssymb,amsthm}
\usepackage{geometry}

\geometry{margin=2.5cm}

\title{其他避免奇异的方法}
\author{第5章和第11章补充说明}
\date{}

\begin{document}

\maketitle

\section{问题提出}

除了SVD和det($J^T J$)方法外，还有哪些方法可以避免或处理奇异点？

\textbf{其他方法}：
\begin{enumerate}
\item 可操作性椭球（Manipulability Ellipsoid）
\item 零空间运动（Null Space Motion）
\item 梯度投影法（Gradient Projection Method）
\item 任务优先级方法（Task Priority Method）
\item 轨迹规划避障（Trajectory Planning）
\item 自适应阻尼（Adaptive Damping）
\end{enumerate}

\section{方法3：可操作性椭球（Manipulability Ellipsoid）}

\subsection{基本原理}

\textbf{可操作性椭球}：描述机器人在不同方向上运动的"容易程度"。

\textbf{定义}：

对于雅可比矩阵$J(\theta) \in \mathbb{R}^{m \times n}$，可操作性椭球定义为：
\begin{equation}
A = J(\theta)J^T(\theta) \in \mathbb{R}^{m \times m}
\end{equation}

\textbf{椭球的特征}：
\begin{itemize}
\item 椭球的形状由$A$的特征值和特征向量决定
\item 椭球的主轴方向：$A$的特征向量方向
\item 椭球的主轴长度：$\sqrt{\lambda_i}$，其中$\lambda_i$是$A$的特征值
\end{itemize}

\subsection{可操作性度量}

\textbf{度量1：条件数}
\begin{equation}
\mu_1 = \frac{\lambda_{\max}}{\lambda_{\min}} = \frac{\sigma_{\max}^2}{\sigma_{\min}^2}
\end{equation}

其中$\lambda_{\max}$和$\lambda_{\min}$是$A$的最大和最小特征值，$\sigma_{\max}$和$\sigma_{\min}$是$J$的最大和最小奇异值。

\textbf{物理意义}：
\begin{itemize}
\item $\mu_1$接近1：椭球接近圆形，远离奇异点
\item $\mu_1$很大：椭球很扁，接近奇异点
\end{itemize}

\textbf{度量2：最小特征值}
\begin{equation}
\mu_2 = \lambda_{\min} = \sigma_{\min}^2
\end{equation}

\textbf{物理意义}：
\begin{itemize}
\item $\mu_2$大：远离奇异点
\item $\mu_2$小：接近奇异点
\end{itemize}

\textbf{度量3：体积}
\begin{equation}
\mu_3 = \sqrt{\det(A)} = \sqrt{\det(JJ^T)} = \prod_{i=1}^r \sigma_i
\end{equation}

\textbf{物理意义}：
\begin{itemize}
\item $\mu_3$大：椭球体积大，远离奇异点
\item $\mu_3$小：椭球体积小，接近奇异点
\item 在奇异点，$\mu_3 = 0$（椭球退化为线段或点）
\end{itemize}

\subsection{避免奇异的策略}

\textbf{策略1：最大化可操作性}

在轨迹规划中，选择最大化可操作性的路径：
\begin{equation}
\theta^* = \arg\max_{\theta} \mu_3(\theta) = \arg\max_{\theta} \sqrt{\det(J(\theta)J^T(\theta))}
\end{equation}

\textbf{策略2：可操作性约束}

在控制中，添加可操作性约束：
\begin{equation}
\mu_3(\theta) \geq \mu_{\text{threshold}}
\end{equation}

如果可操作性低于阈值，调整控制策略。

\textbf{策略3：可操作性加权控制}

在计算伪逆时，根据可操作性调整权重：
\begin{equation}
J^{\dagger} = (J^T J + \lambda(\mu_3) I)^{-1}J^T
\end{equation}

其中$\lambda(\mu_3)$是根据可操作性调整的阻尼系数。

\subsection{优势和劣势}

\textbf{优势}：
\begin{itemize}
\item 提供直观的几何解释
\item 可以量化奇异性
\item 可以用于轨迹规划
\end{itemize}

\textbf{劣势}：
\begin{itemize}
\item 需要计算特征值或奇异值
\item 计算复杂度较高
\item 主要用于离线规划或实时监控
\end{itemize}

\section{方法4：零空间运动（Null Space Motion）}

\subsection{基本原理}

\textbf{零空间}：对于冗余机器人（$n > m$），存在零空间运动。

\textbf{定义}：

雅可比矩阵$J(\theta) \in \mathbb{R}^{m \times n}$的零空间定义为：
\begin{equation}
\mathcal{N}(J) = \{\dot{\theta} : J(\theta)\dot{\theta} = 0\}
\end{equation}

\textbf{性质}：
\begin{itemize}
\item 零空间中的运动不改变末端执行器的速度
\item 零空间的维度：$\dim(\mathcal{N}(J)) = n - \text{rank}(J)$
\item 对于冗余机器人，零空间维度 > 0
\end{itemize}

\subsection{零空间投影矩阵}

\textbf{投影矩阵}：
\begin{equation}
N = I - J^{\dagger}J
\end{equation}

\textbf{性质}：
\begin{itemize}
\item $N$是投影矩阵：$N^2 = N$
\item $JN = 0$：零空间中的运动不影响末端执行器
\item 任何向量可以分解为：$\dot{\theta} = J^{\dagger}V + N\dot{\theta}_{\text{null}}$
\end{itemize}

\subsection{避免奇异的策略}

\textbf{策略1：零空间优化}

在零空间中执行优化运动，避免奇异配置：
\begin{equation}
\dot{\theta} = J^{\dagger}V + N\dot{\theta}_{\text{null}}
\end{equation}

其中$\dot{\theta}_{\text{null}}$是零空间速度，用于优化某些目标（如避免奇异）。

\textbf{零空间速度设计}：
\begin{equation}
\dot{\theta}_{\text{null}} = k \nabla_{\theta} \mu_3(\theta)
\end{equation}

其中$k > 0$是增益，$\nabla_{\theta} \mu_3(\theta)$是可操作性的梯度。

\textbf{物理意义}：
\begin{itemize}
\item 主任务：$J^{\dagger}V$（实现期望的末端执行器速度）
\item 次要任务：$N\dot{\theta}_{\text{null}}$（在零空间中优化，不影响主任务）
\item 次要任务可以是：避免奇异、避免关节限制、优化姿态等
\end{itemize}

\subsection{完整控制律}

\textbf{扩展的逆运动学}：
\begin{equation}
\dot{\theta} = J^{\dagger}V + (I - J^{\dagger}J)\dot{\theta}_{\text{null}}
\end{equation}

\textbf{零空间速度设计}（避免奇异）：
\begin{equation}
\dot{\theta}_{\text{null}} = k \nabla_{\theta} \mu_3(\theta) = k \nabla_{\theta} \sqrt{\det(JJ^T)}
\end{equation}

\textbf{梯度计算}：
\begin{equation}
\nabla_{\theta} \mu_3 = \frac{1}{2\mu_3} \nabla_{\theta} \det(JJ^T)
\end{equation}

\subsection{优势和劣势}

\textbf{优势}：
\begin{itemize}
\item 只适用于冗余机器人
\item 可以在不影响主任务的情况下避免奇异
\item 可以同时优化多个目标
\end{itemize}

\textbf{劣势}：
\begin{itemize}
\item 需要冗余自由度
\item 需要计算可操作性的梯度
\item 计算复杂度较高
\end{itemize}

\section{方法5：梯度投影法（Gradient Projection Method）}

\subsection{基本原理}

\textbf{梯度投影法}：在零空间中执行梯度优化，避免奇异配置。

\textbf{基本公式}：
\begin{equation}
\dot{\theta} = J^{\dagger}V + \alpha N \nabla_{\theta} w(\theta)
\end{equation}

其中：
\begin{itemize}
\item $J^{\dagger}V$：主任务（实现期望速度）
\item $N = I - J^{\dagger}J$：零空间投影矩阵
\item $\nabla_{\theta} w(\theta)$：优化目标的梯度
\item $\alpha$：增益系数
\item $w(\theta)$：优化目标函数（如可操作性）
\end{itemize}

\subsection{优化目标函数}

\textbf{目标1：最大化可操作性}
\begin{equation}
w(\theta) = \mu_3(\theta) = \sqrt{\det(JJ^T)}
\end{equation}

\textbf{目标2：最小化到奇异点的距离}
\begin{equation}
w(\theta) = -\frac{1}{\det(J^T J) + \epsilon}
\end{equation}

其中$\epsilon$是小的正数，避免分母为零。

\textbf{目标3：最大化最小奇异值}
\begin{equation}
w(\theta) = \sigma_{\min}(\theta)
\end{equation}

\subsection{自适应增益}

\textbf{自适应增益}：

根据可操作性调整增益$\alpha$：
\begin{equation}
\alpha = \begin{cases}
\alpha_0 & \text{如果 } \mu_3 > \mu_{\text{threshold}} \\
\alpha_0 \frac{\mu_3}{\mu_{\text{threshold}}} & \text{如果 } \mu_3 \leq \mu_{\text{threshold}}
\end{cases}
\end{equation}

\textbf{物理意义}：
\begin{itemize}
\item 远离奇异点时，使用正常增益
\item 接近奇异点时，增加增益，更积极地避免奇异
\end{itemize}

\subsection{优势和劣势}

\textbf{优势}：
\begin{itemize}
\item 适用于冗余机器人
\item 可以同时优化多个目标
\item 自适应调整
\end{itemize}

\textbf{劣势}：
\begin{itemize}
\item 需要冗余自由度
\item 需要计算梯度
\item 可能影响主任务的性能
\end{itemize}

\section{方法6：任务优先级方法（Task Priority Method）}

\subsection{基本原理}

\textbf{任务优先级}：将任务分为主要任务和次要任务，按优先级执行。

\textbf{主要任务}：末端执行器运动控制
\begin{equation}
V = J(\theta)\dot{\theta}
\end{equation}

\textbf{次要任务}：避免奇异配置
\begin{equation}
w(\theta) = \mu_3(\theta) \geq \mu_{\text{threshold}}
\end{equation}

\subsection{分层控制}

\textbf{第一层：主要任务}
\begin{equation}
\dot{\theta}_1 = J^{\dagger}V
\end{equation}

\textbf{第二层：次要任务（在零空间中）}
\begin{equation}
\dot{\theta}_2 = (I - J^{\dagger}J)\dot{\theta}_{\text{null}}
\end{equation}

\textbf{组合}：
\begin{equation}
\dot{\theta} = \dot{\theta}_1 + \dot{\theta}_2 = J^{\dagger}V + (I - J^{\dagger}J)\dot{\theta}_{\text{null}}
\end{equation}

\subsection{多个次要任务}

\textbf{多个次要任务}：

如果有多个次要任务，可以按优先级处理：
\begin{equation}
\dot{\theta} = J^{\dagger}V + N_1 \dot{\theta}_{\text{null},1} + N_2 \dot{\theta}_{\text{null},2} + \cdots
\end{equation}

其中：
\begin{itemize}
\item $N_1 = I - J^{\dagger}J$：第一级零空间投影
\item $N_2 = I - J_1^{\dagger}J_1$：第二级零空间投影（$J_1$是第一个次要任务的雅可比）
\item 依此类推
\end{itemize}

\subsection{优势和劣势}

\textbf{优势}：
\begin{itemize}
\item 可以处理多个任务
\item 任务优先级明确
\item 适用于复杂任务
\end{itemize}

\textbf{劣势}：
\begin{itemize}
\item 需要冗余自由度
\item 计算复杂
\item 可能影响主任务性能
\end{itemize}

\section{方法7：轨迹规划避障}

\subsection{基本原理}

\textbf{轨迹规划避障}：在轨迹规划阶段就避开奇异配置。

\textbf{策略}：
\begin{enumerate}
\item 识别奇异配置
\item 在配置空间中标记奇异区域
\item 规划轨迹时避开这些区域
\end{enumerate}

\subsection{奇异区域识别}

\textbf{方法1：离线分析}
\begin{itemize}
\item 分析机器人的所有可能配置
\item 计算每个配置的可操作性
\item 标记可操作性低于阈值的区域
\end{itemize}

\textbf{方法2：在线检测}
\begin{itemize}
\item 在轨迹执行过程中实时检测
\item 如果接近奇异点，调整轨迹
\end{itemize}

\subsection{轨迹调整}

\textbf{方法1：路径修改}
\begin{itemize}
\item 检测到接近奇异点时，修改路径
\item 绕过奇异区域
\end{itemize}

\textbf{方法2：速度调整}
\begin{itemize}
\item 在接近奇异点时，降低速度
\item 快速通过奇异区域
\end{itemize}

\subsection{优势和劣势}

\textbf{优势}：
\begin{itemize}
\item 从根本上避免奇异
\item 不需要实时计算
\item 可以离线优化
\end{itemize}

\textbf{劣势}：
\begin{itemize}
\item 需要预先知道奇异配置
\item 可能限制工作空间
\item 轨迹可能不是最优的
\end{itemize}

\section{方法8：自适应阻尼（Adaptive Damping）}

\subsection{基本原理}

\textbf{自适应阻尼}：根据奇异性动态调整阻尼系数。

\textbf{阻尼伪逆}：
\begin{equation}
J^{\dagger} = (J^T J + \lambda(\theta) I)^{-1}J^T
\end{equation}

\textbf{自适应阻尼函数}：

\textbf{方法1：基于det($J^T J$)}
\begin{equation}
\lambda(\theta) = \begin{cases}
0 & \text{如果 } \det(J^T J) > d_{\text{threshold}} \\
\lambda_0\left(1 - \frac{\det(J^T J)}{d_{\text{threshold}}}\right) & \text{如果 } \det(J^T J) \leq d_{\text{threshold}}
\end{cases}
\end{equation}

\textbf{方法2：基于最小奇异值}
\begin{equation}
\lambda(\theta) = \begin{cases}
0 & \text{如果 } \sigma_{\min} > \sigma_{\text{threshold}} \\
\lambda_0\left(1 - \frac{\sigma_{\min}}{\sigma_{\text{threshold}}}\right)^2 & \text{如果 } \sigma_{\min} \leq \sigma_{\text{threshold}}
\end{cases}
\end{equation}

\textbf{方法3：基于可操作性}
\begin{equation}
\lambda(\theta) = \begin{cases}
0 & \text{如果 } \mu_3 > \mu_{\text{threshold}} \\
\lambda_0\left(1 - \frac{\mu_3}{\mu_{\text{threshold}}}\right) & \text{如果 } \mu_3 \leq \mu_{\text{threshold}}
\end{cases}
\end{equation}

\subsection{优势和劣势}

\textbf{优势}：
\begin{itemize}
\item 计算简单
\item 实时性好
\item 平滑过渡
\end{itemize}

\textbf{劣势}：
\begin{itemize}
\item 可能影响控制精度
\item 需要选择合适的阈值
\end{itemize}

\section{方法对比总结}

\subsection{方法对比表}

\begin{center}
\begin{tabular}{|l|l|l|l|}
\hline
\textbf{方法} & \textbf{适用场景} & \textbf{计算复杂度} & \textbf{主要特点} \\
\hline
det($J^T J$) & 所有机器人 & $O(n^3)$ & 检测奇异性 \\
\hline
SVD & 所有机器人 & $O(\min(mn^2, m^2n))$ & 计算稳定伪逆 \\
\hline
可操作性椭球 & 所有机器人 & $O(n^3)$ & 量化奇异性 \\
\hline
零空间运动 & 冗余机器人 & $O(n^3)$ & 不影响主任务 \\
\hline
梯度投影法 & 冗余机器人 & $O(n^3)$ & 优化多个目标 \\
\hline
任务优先级 & 冗余机器人 & $O(n^3)$ & 多任务处理 \\
\hline
轨迹规划 & 所有机器人 & 离线 & 预先避免 \\
\hline
自适应阻尼 & 所有机器人 & $O(n^3)$ & 实时调整 \\
\hline
\end{tabular}
\end{center}

\subsection{选择建议}

\textbf{非冗余机器人（$n = 6$）}：
\begin{itemize}
\item 推荐：SVD + 自适应阻尼
\item 备选：det($J^T J$) + 自适应阻尼
\item 离线：轨迹规划避障
\end{itemize}

\textbf{冗余机器人（$n > 6$）}：
\begin{itemize}
\item 推荐：SVD + 零空间运动
\item 备选：梯度投影法
\item 复杂任务：任务优先级方法
\end{itemize}

\textbf{实时性要求高}：
\begin{itemize}
\item 推荐：det($J^T J$) + 自适应阻尼
\item 备选：自适应阻尼（基于可操作性）
\end{itemize}

\textbf{精度要求高}：
\begin{itemize}
\item 推荐：SVD + 自适应阻尼
\item 备选：零空间运动（冗余机器人）
\end{itemize}

\section{实际应用示例}

\subsection{示例1：7自由度机器人（冗余）}

\textbf{方法组合}：
\begin{enumerate}
\item 使用SVD计算伪逆（数值稳定）
\item 使用零空间运动避免奇异
\item 使用可操作性作为优化目标
\end{enumerate}

\textbf{控制律}：
\begin{equation}
\dot{\theta} = J^{\dagger}V + \alpha(I - J^{\dagger}J)\nabla_{\theta}\mu_3(\theta)
\end{equation}

其中：
\begin{itemize}
\item $J^{\dagger}$：使用SVD计算的阻尼伪逆
\item $\alpha$：根据$\mu_3$自适应调整
\item $\nabla_{\theta}\mu_3(\theta)$：可操作性的梯度
\end{itemize}

\subsection{示例2：6自由度机器人（非冗余）}

\textbf{方法组合}：
\begin{enumerate}
\item 使用det($J^T J$)快速检测
\item 使用自适应阻尼
\item 离线轨迹规划避开已知奇异点
\end{enumerate}

\textbf{控制律}：
\begin{equation}
\dot{\theta} = (J^T J + \lambda(\theta) I)^{-1}J^T V
\end{equation}

其中：
\begin{equation}
\lambda(\theta) = \begin{cases}
0 & \text{如果 } \det(J^T J) > d_{\text{threshold}} \\
\lambda_0\left(1 - \frac{\det(J^T J)}{d_{\text{threshold}}}\right) & \text{否则}
\end{cases}
\end{equation}

\section{总结}

\subsection{关键要点}

\begin{enumerate}
\item \textbf{检测方法}：
\begin{itemize}
\item det($J^T J$)：快速检测
\item 可操作性椭球：量化奇异性
\end{itemize}

\item \textbf{计算方法}：
\begin{itemize}
\item SVD：数值稳定的伪逆
\item 自适应阻尼：实时调整
\end{itemize}

\item \textbf{优化方法}（冗余机器人）：
\begin{itemize}
\item 零空间运动：不影响主任务
\item 梯度投影法：优化多个目标
\item 任务优先级：多任务处理
\end{itemize}

\item \textbf{规划方法}：
\begin{itemize}
\item 轨迹规划：预先避免
\end{itemize}
\end{enumerate}

\subsection{推荐方案}

\textbf{通用方案}（适用于所有机器人）：
\begin{enumerate}
\item 使用det($J^T J$)或可操作性快速检测
\item 使用SVD计算阻尼伪逆
\item 使用自适应阻尼实时调整
\item 离线轨迹规划避开已知奇异点
\end{enumerate}

\textbf{冗余机器人增强方案}：
\begin{enumerate}
\item 上述通用方案
\item 添加零空间运动优化可操作性
\item 使用梯度投影法避免奇异
\end{enumerate}

\end{document}

